\label{sec:discussion}

\subsection{Cross-study synthesis}

The three studies share one methodological skeleton: take the strongest
published quantitative statement available, recompute it from primary data or
a primary model, then push one step past it --- and report both steps with
equal honesty. The push-past step succeeded in Study~B (basin size, learning
curve) and Study~C (benchmark margin, dosage-decodability) and failed
instructively in Study~A (CNN surrogate, retracted). That 2-out-of-3 rate is
not a bug in the programme; it is the expected yield of attempting genuine
extensions rather than rehearsing known results. A programme in which every
extension succeeds is a programme whose extensions were too timid.

\subsection{What the numbers do and do not say}

\textbf{Study A.} The upshift (median $31\%$ above the linear mode, peaking
at intermediate distance from onset) is a property of the Gierer--Meinhardt
kinetics on lattices of the tested size. We have not shown it for other
kinetics (Schnakenberg, FitzHugh--Nagumo), other boundary conditions, or
continuum-accurate solvers; each is a natural next computation and the
pipeline accepts them by parameter change. The non-monotone dependence on
$\varepsilon$ suggests the deviation is controlled by the competition between
mode-coupling strength (which grows with amplitude, i.e.\ with distance from
onset) and band flatness (which shrinks it again), but we have measured, not
derived, this --- a weakly nonlinear analysis in the spirit of amplitude
equations is the obvious theoretical follow-up, and our retracted CNN attempt
is evidence that the effect is not trivially readable from linear data.

\textbf{Study B.} The $0.77\%$ wild-type basin reframes the robustness
question for the segment-polarity module: the model does not pattern
\emph{whatever} the initial state, it patterns given the initial states that
the preceding developmental stages actually produce. The learnability result
($\sim 60$ labelled states suffice) shows the basin boundary, though tiny, is
geometrically simple --- consistent with the ordered dynamics the Derrida
analysis predicts. The extension that matters biologically is the
multi-cellular one, and it is precisely the one that is currently broken
(Section~\ref{sec:negatives}); fixing it is the top open engineering item.

\textbf{Study C.} The CF-decoder margin ($+2.57\%$ EL, CI $[1.99, 3.11]$)
is large relative to the biology: the whole embryo-to-embryo spread of the
furrow is $10.5\%$ EL. Two readings are possible and we cannot yet separate
them: the threshold model leaves profile-shape information unused (our
preferred reading, supported by the ridge model capturing most of the gain
from two shape features), or the CNN exploits line-specific systematics that
happen to survive line-disjoint cross-validation. A second, independent
CF-annotated dataset would decide between them; we flag this as the most
important external validation of the project. The DCLS refinement
(anterior-concentrated dosage coupling) is compatible with the published
supplement figure and adds a spatial statement to it; the
dosage--decodability result extends the same axis to the downstream genes.

\subsection{Open work, in priority order}

\begin{enumerate}
\item Repair or replace the multi-cell segment-polarity extension; validate
against the published two-cell and four-cell patterns.
\item Independent-dataset validation of the CF decoder.
\item Weakly nonlinear theory of the wavenumber upshift; extension to other
kinetics and boundary conditions.
\item Stochastic-target surrogates for pattern selection (predict the
distribution over lattice modes rather than a single wavenumber).
\item Joint decoding across all measured gap genes on the immunofluorescence
panel, approaching the four-gene benchmark of \cite{petkova2019} with the
same honest cross-validation.
\end{enumerate}
